\chapter{Implementation Summary}

\section{Authentication And Sessions}

Authentication is implemented with JWT access tokens and refresh tokens. Passwords are hashed using
\texttt{pbkdf2\_sha256}. Refresh tokens are hashed before storage, and access tokens are tied to the
active refresh-token session so that refresh invalidates older access-token sessions immediately. The
frontend persists the session across page refreshes and refreshes tokens before expiry where necessary.

The local forgot-password flow is intentionally coursework-friendly: instead of integrating an email
provider, the system issues a one-time reset token directly for local evaluation. This preserves the reset
lifecycle without adding external service complexity.

\section{Documents, Roles, And Versioning}

Documents are stored with owner metadata, content, revision identifiers, and version snapshots. Owners can
share directly with other users as editors or viewers, and can also create revocable share-links. All role
enforcement is server-side. The frontend does not merely hide unavailable buttons; the backend rejects
unauthorised requests.

Editing is autosaved and versioned. Revert is implemented as a first-class operation with version history
preview and restoration. Save, share-link creation, and revert support idempotent request identifiers to
avoid duplicate work under retries.

\section{Realtime Collaboration}

Realtime collaboration uses Yjs over a dedicated websocket relay. Clients request a signed collaboration
session from FastAPI, then authenticate to the relay using a websocket subprotocol token. This avoids
placing credentials in a query string while keeping the collaboration layer isolated from the REST API.

The relay handles initial sync, ongoing updates, awareness/presence messages, reconnect behaviour, and
server-side write rejection for viewers. The UI renders active collaborators and remote cursor/selection
state, which also contributes to the bonus criteria.
